% `tex.web` §1090 `init_math`: a display belongs to a *paragraph*.
%
%     begin get_token; {|get_x_token| would fail on \ifmmode!}
%     if (cur_cmd=math_shift)and(mode>0) then @<Go into display math mode@>
%     else  begin back_input; @<Go into ordinary math mode@>;
%       end;
%
% `mode>0` is the whole of it. A restricted horizontal list — the inside of an
% \hbox — is `-hmode`, so `$$` there is not a display: the second `$` is put
% back and dispatched again, and what the box gets is two *empty inline*
% formulas with whatever lies between them set in the box's own font.
%
% We read the second `$` and went into display math wherever we found it, so
% `\hbox{A$$B$$C}` broke a paragraph that did not exist. The box came out
% 115.58682pt wide against TeX's 21.8056, with an \abovedisplayshortskip and a
% centred line inside it, and `B` set from family 1 instead of \rm — every one
% of which is visible on the page, not merely in the transcript.
%
% The trace is the other half: TeX shows `{restricted horizontal mode: math
% shift character $}` and then `{math mode: math shift character $}`, two
% commands for the two characters, because `back_input` sends the second one
% round again. A display shows one.
\catcode`\{=1 \catcode`\}=2 \catcode`\$=3
\tracingonline=1 \tracingcommands=2 \showboxbreadth=99 \showboxdepth=99
\hsize=100pt \parindent=0pt
\font\rm=cmr10 \textfont0=\rm \scriptfont0=\rm \scriptscriptfont0=\rm \rm
\font\mi=cmmi10 \textfont1=\mi \scriptfont1=\mi \scriptscriptfont1=\mi
\font\sy=cmsy10 \textfont2=\sy \scriptfont2=\sy \scriptscriptfont2=\sy
\font\ex=cmex10 \textfont3=\ex \scriptfont3=\ex \scriptscriptfont3=\ex

\setbox0\hbox{A$$B$$C}
\showbox0

% ...and the same inside a \vbox, where the paragraph that `$$` interrupts is
% a real one and the display is real too. Restricted is about the *horizontal*
% list, not about being inside a box.
\setbox1\vbox{\noindent A$$B$$C\par}
\showbox1

% A display does start a paragraph from vertical mode — §1090's `vmode+
% math_shift` backs the `$` up and calls `new_graf(true)`, so the paragraph
% that gets broken is an empty indented one.
\setbox2\vbox{$$B$$ \indent A$X$}
\showbox2
\end
